\documentclass[a4paper]{article}

\usepackage[margin=2cm]{geometry}
\usepackage{amssymb}
\usepackage{pdfpages}
\usepackage{graphbox}
\usepackage{tikz}
\begin{document}

\begin{center}
  {\Large\bf Values of the standard normal distribution function}
\end{center}

\rule{0mm}{5mm}
\pagestyle{empty}
\thispagestyle{empty}
\begin{table}[ht]
\centering
\begin{tabular}{r|llllllllll}
  \hline \rule{0mm}{4mm}
  & {\large .00} & {\large .01} & {\large .02} & {\large .03} & {\large .04}
  & {\large .05} & {\large .06} & {\large .07} & {\large .08} & {\large .09}
  \vphantom{$y_{y_{y_y}}$} \\  \hline
\input{gaussian_table_values.txt}
  \hline
\end{tabular}
\end{table}

\noindent
\(\displaystyle P(Z\leqslant z)=\Phi(z)=\frac{1}{\sqrt{2\pi}}\int_{-\infty}^z\exp\left(-u^2/2\right)\,\mathrm{d}u\)\hfill\hbox{\includegraphics[align=c,width=2.2in]{gaussian_diagram1.pdf}}

\noindent For example, to calculate \(P(Z\leqslant 2.93)\), the last
line gives values for \(2.90-2.99\).  To find the value for 2.93,
choose the entry with column heading ``.03'' to obtain 0.9983.\\ \\



\begin{tikzpicture}[remember picture,overlay]
  \node[anchor=south east,
    xshift=-1cm,
    yshift=1cm]
  at (current page.south east){
    {\tt https://github.com/RobinHankin/tables.git}\rule{90mm}{0mm}\includegraphics[width=1in]{AUT-logo-block.pdf}};
\end{tikzpicture}



\end{document}
